\documentclass[10pt]{article}
\usepackage[T1]{fontenc}
\usepackage[utf8]{inputenc}
\usepackage[french]{babel}
\usepackage{amsmath,amssymb,mathrsfs,mathpazo}
\usepackage{geometry}
\geometry{a4paper,margin=2cm}
\usepackage{graphicx,array,tabularx,tabularray,colortbl,multirow,multicol,booktabs,enumitem,pifont,ulem,fancybox,diagbox,cancel}
\usepackage{pstricks,pstricks-add,pst-plot,pst-tree,pst-node,pst-text,pst-func,pst-math,pst-bspline,pst-3dplot,pst-eucl,pst-solides3d}
\let\pscalculate\undefined
\usepackage{pst-bezier}
\usepackage{tikz,pgfplots,tkz-tab}
\usetikzlibrary{arrows,arrows.meta,patterns,calc,positioning,trees,decorations.markings,shapes.geometric,shapes.misc,shapes.arrows,intersections,decorations.pathmorphing,decorations.pathreplacing,backgrounds,shadows,babel}
\pgfplotsset{compat=1.18}
\usepgfplotslibrary{fillbetween}
\usepackage[np]{numprint}
\usepackage{listings,makecell,pas-tableur,twoopt}
\let\barre\undefined
\usepackage[pstricks]{tablvar}
\definecolor{darkgreen}{RGB}{0,100,0}
\providecommand{\red}{\color{red}}
\providecommand{\blue}{\color{blue}}
\providecommand{\green}{\color{green}}
\providecommand{\black}{\color{black}}
\providecommand{\white}{\color{white}}
\providecommand{\gray}{\color{gray}}
\providecommand{\cyan}{\color{cyan}}
\providecommand{\magenta}{\color{magenta}}
\providecommand{\yellow}{\color{yellow}}
\providecommand{\violet}{\color{violet}}
\providecommand{\orange}{\color{orange}}
\providecommand{\euro}{\texteuro}
\providecommand{\EUR}{\texteuro}
\providecommand{\ieme}{\textsuperscript{e}}
\providecommand{\degres}{\ensuremath{^\circ}}
\providecommand{\degree}{\ensuremath{^\circ}}
\providecommand{\vv}[1]{\overrightarrow{#1}}
\providecommand{\Coord}[1]{#1}
\providecommand{\decosix}{}
\providecommand{\starredbullet}{\ensuremath{\bullet}}
\pagestyle{empty}
\setlength{\parindent}{0pt}
\usepackage{adjustbox}\DeclareUnicodeCharacter{2236}{:}\DeclareUnicodeCharacter{2190}{\ensuremath{\leftarrow}}
\providecommand{\R}{\mathbb{R}}
\providecommand{\N}{\mathbb{N}}
\providecommand{\Z}{\mathbb{Z}}
\providecommand{\Q}{\mathbb{Q}}
\providecommand{\D}{\mathbb{D}}
\providecommand{\C}{\mathbb{C}}
\providecommand{\e}{{\mathrm{e}}}
\providecommand{\vect}[1]{\overrightarrow{#1}}
\providecommand{\vectt}[1]{\overrightarrow{\mathrm{#1}}}
\providecommand{\vvt}[1]{\overrightarrow{\mathrm{#1}}}
\providecommand{\np}[2][]{{#2~#1}}
\providecommand{\vv}[1]{{\overrightarrow{#1}}}
\renewcommand{\d}{{\mathrm{d}}}
\providecommand{\degres}{{^\circ}}
\providecommand{\degre}{{^\circ}}
\providecommand{\blue}{\color{blue}}
\providecommand{\red}{\color{red}}
\providecommand{\enskip}{\quad}
\providecommand{\llbracket}{\lbrack\!\lbrack}
\providecommand{\rrbracket}{\rbrack\!\rbrack}
\providecommand{\euro}{{\text{€}}}
\providecommand{\no}{{\text{n°}}}
\providecommand{\interval}[2]{{[#1;#2]}}
\providecommand{\intervFO}[2]{{[#1;#2[}}
\providecommand{\dotfill}{\ldots}

\providecommand{\ds}{\displaystyle}
\providecommand{\textding}[1]{#1}
\providecommand{\Oij}{\(\mathrm{O};\vec{\imath},\vec{\jmath}\)}
\providecommand{\Oijk}{\(\mathrm{O};\vec{\imath},\vec{\jmath},\vec{k}\)}
\providecommand{\Ouv}{\(\mathrm{O};\vec{u},\vec{v}\)}
\providecommand{\intervFF}[2]{\ensuremath{[#1;#2]}}
\renewcommand{\euro}{\ifmmode\text{\texteuro}\else\texteuro\fi}
\DeclareUnicodeCharacter{00B0}{\ifmmode^{\circ}\else\textdegree\fi}
\DeclareUnicodeCharacter{20AC}{\ifmmode\text{\texteuro}\else\texteuro\fi}
% Énoncés : Éducation nationale. Transcription et corrigés : APMEP, Denis Vergès et contributeurs. Découpage et classement : CoopMaths, Sébastien Lozano, Éric Elter et contributeurs. Adaptation de la présentation : Maths Explorer.
% CC BY-SA : https://coopmaths.fr/www/about/
% https://coopmaths.fr/alea/static/bac/2024/tex/bac_2024_06_sujet2_metropole_1_cor.tex
\begin{document}
\noindent\textit{APMEP / CoopMaths. CC BY-SA. Présentation adaptée par Maths Explorer.}\par\medskip
\medskip

Puisque l'on interroge un étudiant au hasard, on est en situation d'équiprobabilité, et les proportions sont assimilables à des probabilités.

\begin{enumerate}
	\item Comme 91,7\,\% des étudiants ont répondu oui, on a donc $P(Q) = 0,917$.

	Et $P_{\overline{R}}\big( \overline{Q}\big) = 0,65$, cela correspond à la probabilité que le candidat interrogé ait répondu \og non\fg{}, sachant qu'il a échoué à l'examen.

	\item \begin{enumerate}
		\item L'arbre pondéré complété est :

		\begin{center}
			\begin{tikzpicture}[xscale=1,yscale=1,baseline={(R.base)}]
				
				\tikzstyle{fleche}=[-,>=latex,thick]
				\tikzstyle{noeud}=[fill=white,circle,inner sep=2pt]
				\tikzstyle{feuille}=[fill=white,circle,inner sep=2pt]
				\tikzstyle{etiquette}=[pos=0.6,fill=white, inner xsep=3pt, inner ysep=1.5pt]
				
				\def\DistanceInterNiveaux{3}
				\def\DistanceInterFeuilles{0.8}
				
				\def\NiveauA{(0)*\DistanceInterNiveaux}
				\def\NiveauB{(1)*\DistanceInterNiveaux}
				\def\NiveauC{(2)*\DistanceInterNiveaux}
				\def\InterFeuilles{(-1)*\DistanceInterFeuilles}
				
				\node[noeud] (R) at ({\NiveauA},{(1.5)*\InterFeuilles}) {};
				\node[noeud] (Ra) at ({\NiveauB},{(0.5)*\InterFeuilles}) {$R$};
				\node[feuille] (Raa) at ({\NiveauC},{(0)*\InterFeuilles}) {$Q$};
				\node[feuille] (Rab) at ({\NiveauC},{(1)*\InterFeuilles}) {$\overline{Q}$};
				\node[noeud] (Rb) at ({\NiveauB},{(2.5)*\InterFeuilles}) {$\overline{R}$};
				\node[feuille] (Rba) at ({\NiveauC},{(2)*\InterFeuilles}) {$Q$};
				\node[feuille] (Rbb) at ({\NiveauC},{(3)*\InterFeuilles}) {$\overline{Q}$};
				
				\draw[fleche] (R.east)--(Ra.west)node[etiquette] {$x$};
				\draw[fleche] (Ra.east)--(Raa.west) node[etiquette] {$0,98 $};
				\draw[fleche] (Ra.east)--(Rab.west) node[etiquette] {$0,02 $};
				\draw[fleche] (R.east)--(Rb.west) node[etiquette] {$1-x $};
				\draw[fleche] (Rb.east)--(Rba.west) node[etiquette] {$0,35 $};
				\draw[fleche] (Rb.east)--(Rbb.west) node[etiquette] {$0,65 $};
			\end{tikzpicture}
		\end{center}

		\item On sait que $P(Q)=0,917$, or, comme $R$ et $\overline{R}$ partitionnent l'univers, en vertu de la loi des probabilités totales, on a également :

		$\begin{aligned} P(Q) &= P\big(Q \cap R\big) + P\left(Q \cap \overline{R}\right)\\
		&= x\times 0,98 + (1 -x) \times 0,35\\
		&= 0,98x + 0,35 - 0,35x\\
		&= 0,63x + 0,35		\end{aligned}$

		Puisque la probabilité de $Q$ est unique, on peut donc écrire, puis résoudre l'équation :
		\quad
		$\begin{aligned}[t] 0,63x + 0,35 = 0,917 &\iff 0,63x = 0,567\\
		&\iff x = \dfrac{0,567}{0,63}\\
		&\iff x = 0,9\end{aligned}$

La solution de l'équation est donc 0,9 : la probabilité que l'étudiant ait réussi son examen est donc de 0,9 (ou bien : 90\,\% des étudiants ont réussi leur examen).
	\end{enumerate}

	\item Pour cette question, on doit calculer $P_Q(R)$.

	On a :\quad $P_Q(R) = \dfrac{P(Q \cap R)}{P(Q)} = \dfrac{0,9 \times 0,98}{0,917}=\dfrac{0,882}{0,917}=\dfrac{126}{131}\approx0,962$

Si l'étudiant interrogé a répondu \og oui\fg{}  à la question, la probabilité qu'il ait réussi l'examen est d'environ 0,962 (à $10^{-3}$ près).

C'est-à-dire qu'environ 96,2\,\% des étudiants ayant répondu \og oui\fg{} ont effectivement réussi l'examen.

	\item Par exploration à la calculatrice, avec une loi binomiale de paramètres $(20~;~0,615)$, on a :\quad
	$P(N\geqslant 11)\approx 0,797$, \quad et \quad $P(N \geqslant 12) \approx 0,649$.

En choisissant de récompenser les candidats dont la note est supérieure ou égale à 11, la directrice récompensera environ 80\,\% d'eux.

	\item Comme les dix variables aléatoires suivent la même loi binomiale de paramètres 20 et 0,615, chacune de ces variables aléatoires a

	\begin{itemize}[label=\textbullet]
		\item pour espérance : \quad $E(N_i) = n \times p = 20 \times 0,615 = 12,3$;
		\item et pour variance : \quad $V(N_i) = n \times p \times (1-p) = 20 \times 0,615 \times 0,385 = 4,7355$.
	\end{itemize}

$S$ étant la somme de ces dix variables aléatoires, on a :

	\begin{itemize}
		\item $E(S) = E(N_1) + E(N_2) + \dots + E(N_{10}) = 12,3 + 12,3 + ... + 12,3 = 123$;
		\item comme les variables aléatoires sont indépendantes,

		$V(S) =  V(N_1) + V(N_2) + \dots + V(N_{10}) = 10 \times 4,7355 = 47,355$	.
\end{itemize}

	\item
	\begin{enumerate}
		\item  La variable aléatoire $M$ donne la moyenne des notes obtenues par un groupe (un échantillon) de 10 étudiants choisis au hasard.

		\item On a bien
		\begin{itemize}
			\item $E(M)= \dfrac{1}{10}\times E(S) = \dfrac{1}{10}\times 123 = 12,3$;
			\item $V(M)=\left(\dfrac{1}{10}\right)^2\times V(S) = \dfrac{1}{100}\times 47,355 =  \np{0,47355}$.
		\end{itemize}
		\item On remarque que l'intervalle de notes considéré : $]10,3 ~;~ 14,3[$ est un intervalle centré sur l'espérance de $M$, autrement dit, la probabilité qui nous intéresse est celle de l'évènement :
		\quad $A = \big\{|M - E(M)| < 2\big\} $.

		L'évènement contraire de cet évènement est : 		\quad $\overline{A} = \big\{|M - E(M)| \geqslant 2\big\} $.

		L'inégalité de Bienaymé-Tchebychev affirme que :

		$\begin{aligned}[t] P\big( |M - E(M)| \geqslant 2 \big) \leqslant \dfrac{V(M)}{2^2 } &\iff P\big( |M -12,3| \geqslant 2 \big) \leqslant \dfrac{0,47355}{4}\\
		&\iff P\big( |M -12,3| \geqslant 2 \big) \leqslant 0,1183875 \\
		&\iff P\big( |M -12,3| <  2 \big) \geqslant 1 - 0,1183875\\
		&\iff P\big( |M -12,3| <  2 \big) \geqslant 0,8816125  \end{aligned}$

		L'affirmation \og La probabilité que la moyenne des notes de dix étudiants pris au hasard soit strictement comprise entre 10,3 et 14,3 est d'au moins $80\,\%$ \fg{} est donc bien correcte, la probabilité que cet évènement soit réalisé étant supérieure à 0,88, elle est aussi supérieure à 0,8 qui correspond à 80\,\%.
	\end{enumerate}
\end{enumerate}

\medskip

\end{document}
